\subsection{Statistical Analysis and Bias Mitigation}
\label{subsec:statistical-analysis}

A critical threat to internal validity in software engineering experiments is \textit{experimenter bias}, particularly confirmation bias and the learning effect across repeated trials \cite{wohlin_2012}. In studies evaluating diagnostic latency, this threat is magnified when the artifact designer acts as the evaluator, as prior knowledge of fault loci artificially deflates diagnostic latency and undermines trial independence \cite{sjoberg_2007, wang_2024}. Furthermore, as systematized by \cite{barata_2026}, empirical evaluations in distributed microservice architectures demand high diagnostic consistency and reliability, which are compromised by human cognitive variability or non-deterministic reasoning mechanisms.

To eliminate these confounding factors, our experimental design decouples fault generation from root-cause identification by employing an automated, algorithmic Root Cause Analysis (RCA) engine grounded in multimodal microservice diagnosis \cite{wang_2024, barata_2026}. The fault injection module operates as a blind, randomized process: at each iteration, the target microservice and perturbation parameters are selected at runtime without disclosing state variables to the diagnostic engine. This programmatic separation guarantees that diagnostic convergence depends strictly on the information accessibility and topological correlation provided by each architectural baseline.

To ensure objective operationalization, explicit convergence criteria govern both automated diagnostic workflows:

\begin{itemize}
    \item \textbf{Vanilla Algorithmic Workflow (Scenario 3):} Success ($Acc = 1$) is recorded if the automated traversal engine correctly identifies the ground-truth perturbed microservice (e.g., attributing anomalous downstream latency to the target microservice) using exclusively native Kubernetes API telemetry. The engine executes breadth-first log inspections and native \texttt{metrics-server} queries. The diagnostic timestamp ($T_{\text{diag}}$) is recorded programmatically at the completion of the algorithmic decision loop. Failure ($Acc = 0$) occurs if the engine reports an incorrect component or fails to converge within the operational boundary ($T_{\text{max}} = 300\,\text{s}$). In timed-out trials, latency is right-censored at $300\,\text{s}$.
    
    \item \textbf{PANOPTES Algorithmic Workflow (Scenario 4):} Success ($Acc = 1$) is achieved if the diagnostic engine isolates the true faulty microservice by traversing the structured call graph and span dependencies exposed by the Tempo and Prometheus APIs \cite{wang_2024}. The timestamp ($T_{\text{diag}}$) is marked immediately upon isolating the anomalous span path and its correlated contextual log signature. If the engine fails to isolate the target or exceeds $T_{\text{max}} = 300\,\text{s}$, the run is flagged as a failure ($Acc = 0$) and right-censored at $300\,\text{s}$.
\end{itemize}

The empirical datasets gathered across all $N = 30$ iterations per scenario are evaluated via inferential statistical testing mapped directly to the three research hypotheses:

\begin{enumerate}
    \item \textbf{Diagnostic Accuracy Analysis (Testing $H_{1,2}$):} Binary diagnostic outcomes ($1 = \text{Success}, 0 = \text{Failure}$) are structured into a $2\times2$ contingency table comparing the Vanilla and PANOPTES configurations. We employ \textbf{Fisher's Exact Test} (one-tailed, $\alpha = 0.05$). This test calculates the exact hypergeometric probability of the observed distribution, eliminating the approximations of the Chi-Square test under sparse or unbalanced contingency cells \cite{agresti_2007}. A statistically significant result ($p < 0.05$) rejects $H_{0,2}$ in favor of $H_{1,2}$.

    \item \textbf{Diagnostic Latency Analysis (Testing $H_{1,1}$):} To test differences in $MTTDiag$ between Scenario 3 and Scenario 4, we enforce a distribution-dependent decision framework \cite{wohlin_2012}:
    \begin{enumerate}
        \item \textbf{Normality Assessment:} All latency distributions are evaluated using the \textbf{Shapiro-Wilk test} ($\alpha = 0.05$).
        \item \textbf{Inferential Test:} If both distributions satisfy normality, we apply the two-sample \textit{Independent Samples t-test} (or \textit{Welch's t-test} under unequal variances). If normality is violated---the expected distribution profile for right-censored troubleshooting latencies \cite{barata_2026}---we apply the non-parametric \textbf{Mann-Whitney U Test} (two-tailed, $\alpha = 0.05$). A statistically significant rank-sum reduction ($p < 0.05$) rejecting $H_{0,1}$ confirms the diagnostic efficiency hypothesis ($H_{1,1}$).
    \end{enumerate}

    \item \textbf{Bounded Overhead Compliance (Testing $H_{1,3}$):} To evaluate whether the static footprint of the observability stack respects the operational envelope, the steady-state aggregate overhead ($R_{\text{PANOPTES}}$ from Scenario 2) is evaluated against the bounded threshold $\theta = 0.15 \cdot \sum R_{\text{workload}}$:
    \begin{enumerate}
        \item The distribution of steady-state CPU (millicores) and RAM (RSS MiB) samples is evaluated for normality via the \textbf{Shapiro-Wilk test}.
        \item We apply a one-tailed \textbf{One-Sample Wilcoxon Signed-Rank Test} (or a \textit{One-Sample t-test} if normally distributed) against the constant boundary $\theta$. Rejection of $H_{0,3}$ ($p < 0.05$) demonstrates that the operational overhead of PANOPTES remains statistically bounded within acceptable operational limits ($H_{1,3}$).
    \end{enumerate}
\end{enumerate}

By combining an automated deterministic RCA engine, blind randomized fault injection, right-censoring for unresolved trials, and non-parametric hypothesis tests, this evaluation protocol eliminates experimenter bias and guarantees strict mathematical replicability.